\documentclass[12pt,letterpaper]{article}

\usepackage[utf8]{inputenc}
\usepackage[spanish]{babel}
\usepackage{amsmath}
\usepackage{amssymb}
\usepackage{geometry}
\usepackage{listings}
\usepackage{xcolor}
\usepackage{hyperref}
\usepackage{booktabs}
\usepackage{graphicx}
\usepackage{lmodern}
\usepackage{float}
\geometry{left=2.5cm,right=2.5cm,top=2.5cm,bottom=2.5cm}

% Configuración para bloques de código
\definecolor{codegreen}{rgb}{0,0.6,0}
\definecolor{codegray}{rgb}{0.5,0.5,0.5}
\definecolor{codepurple}{rgb}{0.58,0,0.82}
\definecolor{backcolour}{rgb}{0.95,0.95,0.92}

\lstdefinestyle{mystyle}{
    backgroundcolor=\color{backcolour},   
    commentstyle=\color{codegreen},
    keywordstyle=\color{magenta},
    numberstyle=\tiny\color{codegray},
    stringstyle=\color{codepurple},
    basicstyle=\ttfamily\footnotesize,
    breakatwhitespace=false,         
    breaklines=true,                 
    captionpos=b,                    
    keepspaces=true,                 
    numbers=left,                    
    numbersep=5pt,                  
    showspaces=false,                
    showstringspaces=false,
    showtabs=false,                  
    tabsize=2
}
\lstset{style=mystyle}

\title{\textbf{Desarrollo del Gemelo Digital del Beamline: Implementación de la Fase 2 (Emulador Forward)}}
\author{Proyecto MIT - UNAL, Medellín \\ Hackathon Digital Twin}
\date{\today}

\begin{document}

\maketitle


\section{Introducción}
Este documento describe el desarrollo y la implementación en código de la \textbf{Fase 2: El Emulador Forward (DNN)} en el marco del Gemelo Digital del beamline de iones. La arquitectura sigue la filosofía del paper \textit{CBOL-Tuner} de Rautela et al., donde el modelado predictivo del sistema no se realiza mapeando directamente las acciones físicas a los estados complejos (de muy alta dimensión y ruidosos), sino reduciendo el mapeo a través del espacio latente entrenado en la Fase 1.

A continuación, se justifica la fundamentación matemática, la arquitectura de la red implementada, el proceso de entrenamiento y la verificación práctica de la diferenciabilidad del emulador.

\section{Fundamentación Teórica del Emulador Forward}
En un beamline real o simulado mediante SIMION, el comportamiento del sistema está determinado por un vector de voltajes de entrada (acciones $y \in \mathbb{R}^8$) y produce una distribución de partículas en el detector (estado $X \in \mathbb{R}^{400}$).

\subsection{¿Qué es el Emulador Forward?}
El Emulador Forward es un sistema de red neuronal compuesto por dos módulos paralelos diferenciables que reciben el vector de voltajes de entrada $y \in \mathbb{R}^8$:
\begin{enumerate}
    \item \textbf{Forward Regressor ($f_\psi$):} Mapea los voltajes a las coordenadas del espacio latente de la forma:
    \begin{equation}
        f_\psi: y \mapsto \hat{z} \in \mathbb{R}^4
    \end{equation}
    \item \textbf{Transmission Regressor ($R_\omega$):} Mapea los voltajes a la transmisión fraccionaria del haz:
    \begin{equation}
        R_\omega: y \mapsto T \in [0, 1]
    \end{equation}
\end{enumerate}

El flujo completo de predicción en el emulador compuesto (\texttt{FullEmulator}) para reconstruir el perfil espacial escalado en el detector es:
\begin{equation}
    \hat{X} = R_\omega(y) \times \text{Decoder}(f_\psi(y))
\end{equation}
donde $\text{Decoder}(\hat{z})$ produce la distribución espacial de probabilidad (forma) y $R_\omega(y)$ escala la intensidad según la corriente transmitida.

\subsection{¿Por qué mapear al Espacio Latente en lugar de mapear directo?}
Mapear directamente los voltajes a la distribución del detector ($y \to X$) presenta múltiples inconvenientes:
\begin{itemize}
    \item \textbf{Alta dimensionalidad y escasez de datos:} Mapear 8 entradas a 400 salidas independientes requiere una gran cantidad de parámetros en la red, lo que induce sobreajuste con pocos datos.
    \item \textbf{Predicciones físicamente imposibles:} Una red que mapee directamente a $X$ no conoce las correlaciones espaciales reales y puede predecir ruido de fondo o distribuciones no físicas (hallucinations).
    \item \textbf{Restricción física (Manifold Latente):} Al forzar a que el regresor apunte a la zona latente $z$ del VAE, las predicciones del decodificador están restringidas al subespacio (manifold) de distribuciones que el VAE aprendió como ``físicamente posibles'' durante la Fase 1.
\end{itemize}

\subsection{¿Para qué sirve esta arquitectura? (Rúbrica D)}
El fin de este diseño es la \textbf{Diferenciabilidad Analítica de Extremo a Extremo}. Como los regresores ($f_\psi$, $R_\omega$) y el Decodificador del VAE están implementados en PyTorch, el flujo completo:
\begin{equation}
    \hat{X} = R_\omega(y) \times \text{Decoder}(f_\psi(y))
\end{equation}
se compone enteramente de operaciones diferenciables. Esto nos permite calcular de manera analítica e instantánea los gradientes del perfil predicho con respecto a los voltajes físicos de entrada $y$:
\begin{equation}
    \nabla_y \hat{X} = \frac{\partial \hat{X}}{\partial y}
\end{equation}
usando retropropagación (\texttt{.backward()}). Gracias a esta diferenciabilidad, es posible optimizar de forma determinista el voltaje de los electrodos en milisegundos para maximizar la transmisión escalar o corregir la órbita geométrica del haz, evitando realizar simulaciones costosas a ciegas.

\section{Arquitectura Neuronal del Emulador}
La red se implementa en PyTorch en el archivo \texttt{model.py} y se compone de dos elementos fundamentales:

\subsection{Forward Regressor y Transmission Regressor}
El emulador consta de dos arquitecturas MLP paralelas:
\begin{itemize}
    \item \textbf{Forward Regressor:} Mapea $\mathbb{R}^8 \to \mathbb{R}^4$ (espacio latente) mediante:
    \begin{equation}
        8 \xrightarrow{\text{Linear+ReLU}} 64 \xrightarrow{\text{Linear+ReLU}} 128 \xrightarrow{\text{Linear+ReLU}} 64 \xrightarrow{\text{Linear}} 4
    \end{equation}
    \item \textbf{Transmission Regressor:} Mapea $\mathbb{R}^8 \to [0, 1]$ (transmisión) mediante:
    \begin{equation}
        8 \xrightarrow{\text{Linear+ReLU}} 64 \xrightarrow{\text{Linear+ReLU}} 128 \xrightarrow{\text{Linear+ReLU}} 64 \xrightarrow{\text{Linear}} 1 \xrightarrow{\text{Sigmoid}} T
    \end{equation}
\end{itemize}
No se incluye Batch Normalization en estas redes para garantizar que el comportamiento sea determinista y suave respecto a los voltajes físicos de entrada.

\subsection{Full Emulator (Compuesto)}
La clase \texttt{FullEmulator} amarra ambas redes y el Decodificador del VAE en un único pipeline de PyTorch. Al instanciarse, congela los pesos del decodificador:
\begin{lstlisting}[language=Python]
for param in self.decoder.parameters():
    param.requires_grad = False
\end{lstlisting}
Esto garantiza que durante el cálculo de gradientes para control, no se altere la estructura generativa del espacio latente.

\section{Proceso de Entrenamiento y Datos}
El entrenamiento está gestionado por el archivo \texttt{train\_dnn.py}.

\subsection{Entrenamiento de Redes y Balanceo de Clases}
\begin{itemize}
    \item \textbf{Forward Regressor:} Se entrena únicamente en las 41 muestras exitosas ($T > 0$), ya que el espacio latente de formas solo tiene representación física válida para haces que llegan al detector. Los histogramas reales se pasan por el \texttt{Encoder} para construir los objetivos $z_i = \mu_{z,i} \in \mathbb{R}^4$. Se optimiza mediante la pérdida MSE:
    \begin{equation}
        \mathcal{L}_{\text{MSE\_z}} = \frac{1}{B} \sum_{b=1}^{B} \|\hat{z}_b - z_b\|^2
    \end{equation}
    \item \textbf{Transmission Regressor:} Se entrena sobre las 486 muestras del dataset. Para mitigar el colapso debido al 92\% de muestras con transmisión nula, se realiza un oversampling de las muestras positivas en el loader de PyTorch. Se optimiza mediante pérdida MSE contra el valor real de transmisión:
    \begin{equation}
        \mathcal{L}_{\text{MSE\_T}} = \frac{1}{B} \sum_{b=1}^{B} (\hat{T}_b - T_b)^2
    \end{equation}
\end{itemize}
Ambos modelos se optimizan usando Adam con una tasa de aprendizaje de $10^{-3}$.

\section{Base de Código de la Fase 2}
Los códigos correspondientes se ubican en la carpeta \texttt{src/phase\_2\_DNN/}:
\begin{itemize}
    \item \textbf{\texttt{model.py}}: Contiene las clases \texttt{ForwardRegressor}, \texttt{TransmissionRegressor} y el emulador compuesto \texttt{FullEmulator}.
    \item \textbf{\texttt{train\_dnn.py}}: Carga el dataset, el VAE de la Fase 1, realiza la codificación latente para muestras positivas, aplica el oversampling, entrena ambos regresores y guarda sus pesos correspondientes en \texttt{forward\_regressor.pt} y \texttt{transmission\_regressor.pt}.
    \item \textbf{\texttt{verify\_emulator.py}}: Script de diagnóstico que verifica la diferenciabilidad calculando los gradientes de la salida del emulador compuesto respecto a los voltajes físicos.
\end{itemize}

\section{Resultados y Verificación}
El regresor forward se entrenó sobre las muestras exitosas codificadas en 4D, y el regresor de transmisión sobre todo el dataset balanceado.

\subsection{Curva de Convergencia}
El entrenamiento de 250 épocas en CUDA arrojó los siguientes resultados de validación óptimos:
\begin{itemize}
    \item \textbf{Forward Regressor ($y \to z \in \mathbb{R}^4$):} MSE de validación óptimo de \textbf{1.249795}. La consistencia en las 4 dimensiones latentes es sumamente robusta.
    \item \textbf{Transmission Regressor ($y \to T$):} MSE de validación óptimo de \textbf{0.000213}. La resolución de la frontera física y la corriente transmitida presentan un error despreciable.
\end{itemize}
Estos resultados confirman que el emulador predice tanto la geometría fina del haz como la fracción de transmisión física con excelente fidelidad.

\subsection{Cálculo Analítico de Sensibilidad (Gradientes)}
Se ejecutó el script de verificación cargando el emulador compuesto. Introduciendo voltajes de prueba normalizados se predijo el perfil de haz $\hat{X}$ y se definió la transmisión como la suma de bins del histograma: $\text{Transmission} = \sum \hat{X}_i$. Aplicando retropropagación, obtuvimos los gradientes analíticos reales en consola:
\begin{table}[H]
    \centering
    \caption{Gradientes de sensibilidad obtenidos mediante retropropagación.}
    \begin{tabular}{lc}
        \toprule
        \textbf{Voltaje del Electrodo} & \textbf{Sensibilidad Analítica ($\partial \text{Transmisión} / \partial V$)} \\
        \midrule
        $V_3$ (Lente Einzel 1 centro) & $-0.019970$ \\
        $V_6$ (Lente Einzel 2 centro) & $-0.023670$ \\
        $V_9$ (Curvador cuadrupolar 1) & $+0.070732$ \\
        $V_{10}$ (Curvador cuadrupolar 2) & $-0.024156$ \\
        $V_{11}$ (Curvador cuadrupolar 3) & $+0.022326$ \\
        $V_{12}$ (Curvador cuadrupolar 4) & $-0.067548$ \\
        $V_{15}$ (Placa dirección 1) & $+0.055127$ \\
        $V_{18}$ (Placa dirección 2) & $+0.001081$ \\
        \bottomrule
    \end{tabular}
\end{table}

Estos gradientes muestran, por ejemplo, que para incrementar la transmisión, se debería aumentar ligeramente el voltaje en el electrodo 9 ($V_9$) y disminuir el de la lente Einzel 1 ($V_3$). Esto valida la completa diferenciabilidad y la utilidad de control del Gemelo Digital.

\end{document}
